% Document/LinearAlgebra/VectorSpaces/Matrices/RankNullityChangeOfBasis.tex
% Phase 16g: Rank/Nullity Invariance and Change of Bases Formula

\newpage
\section{Rank, Nullity, and Change of Bases}
\label{sec:rank-nullity-change-of-bases}

\subsection{Rank and Nullity are Basis-Independent}

\begin{theorem}[Rank and Nullity Invariance]
\label{thm:rank-nullity-invariance}
    Let $T: U \to V$ be $K$-linear with ordered bases $\mathcal{B}_U$, $\mathcal{B}_V$.
    Then:
    \[
        \Rank{T} = \Rank{\CoordMatrix[\mathcal{B}_U][\mathcal{B}_V]{T}},
        \qquad
        \Nullity{T} = \Nullity{\CoordMatrix[\mathcal{B}_U][\mathcal{B}_V]{T}}
    \]
\end{theorem}

\begin{proof}
    The commutative square gives
    $\CoordMatrix[\mathcal{B}_U][\mathcal{B}_V]{T}
    = \Composition{\phi_{\mathcal{B}_V}}{\Composition{T}{\phi_{\mathcal{B}_U}^{-1}}}$.
    Since $\phi_{\mathcal{B}_U}$ and $\phi_{\mathcal{B}_V}$ are isomorphisms, pre- and
    post-composing by them does not change rank or nullity: by
    Theorem~\ref{thm:bijectivity}, isomorphisms preserve the dimension of every subspace.
    Explicitly:
    \begin{itemize}
        \item $\Rng{\CoordMatrix[\mathcal{B}_U][\mathcal{B}_V]{T}}
              = \phi_{\mathcal{B}_V}(\Rng{T})$, so
              $\Rank{\CoordMatrix[\mathcal{B}_U][\mathcal{B}_V]{T}}
              = \Dim{\phi_{\mathcal{B}_V}(\Rng{T})}
              = \Dim{\Rng{T}}
              = \Rank{T}$.
        \item $\Ker{\CoordMatrix[\mathcal{B}_U][\mathcal{B}_V]{T}}
              = \phi_{\mathcal{B}_U}(\Ker{T})$, so
              $\Nullity{\CoordMatrix[\mathcal{B}_U][\mathcal{B}_V]{T}}
              = \Dim{\phi_{\mathcal{B}_U}(\Ker{T})}
              = \Dim{\Ker{T}}
              = \Nullity{T}$. \qedhere
    \end{itemize}
\end{proof}

\begin{remark}
    $\Rank{T}$ and $\Nullity{T}$ can be computed from any matrix representation of $T$;
    the choice of bases is irrelevant.
\end{remark}

\subsection{Change of Bases Formula}

\begin{theorem}[Change of Bases Formula]
\label{thm:change-of-bases}
    Let $T: U \to V$ be $K$-linear, let $\mathcal{B}_U$, $\mathcal{C}_U$ be two ordered
    bases of $U$, and $\mathcal{B}_V$, $\mathcal{C}_V$ two ordered bases of $V$. Then:
    \[
        \CoordMatrix[\mathcal{B}_U][\mathcal{B}_V]{T}
        = \CoordMatrix[\mathcal{C}_V][\mathcal{B}_V]{\Identity{V}}
          \cdot \CoordMatrix[\mathcal{C}_U][\mathcal{C}_V]{T}
          \cdot \CoordMatrix[\mathcal{B}_U][\mathcal{C}_U]{\Identity{U}}
    \]
\end{theorem}

\begin{proof}
    Apply Corollary~\ref{thm:comp-mult} (composition equals multiplication) to the
    trivial factorization $T = \Composition{\Identity{V}}{\Composition{T}{\Identity{U}}}$,
    routing through the intermediate bases $\mathcal{C}_U$ on $U$ and $\mathcal{C}_V$ on $V$:
    \[
        \CoordMatrix[\mathcal{B}_U][\mathcal{B}_V]{T}
        = \CoordMatrix[\mathcal{B}_U][\mathcal{B}_V]
          {\Composition{\Identity{V}}{\Composition{T}{\Identity{U}}}}
        = \CoordMatrix[\mathcal{C}_V][\mathcal{B}_V]{\Identity{V}}
          \cdot \CoordMatrix[\mathcal{C}_U][\mathcal{C}_V]{T}
          \cdot \CoordMatrix[\mathcal{B}_U][\mathcal{C}_U]{\Identity{U}}
        \qedhere
    \]
\end{proof}

\begin{remark}[Commutative Grid]
    The commutative square from Definition~\ref{def:matrix-of-map} extends to a $3 \times 2$
    grid when a second pair of bases is introduced:
    \[
    \begin{tikzcd}[column sep=5em, row sep=3em]
        U \arrow[r, "T"] \arrow[d, "\phi_{\mathcal{B}_U}"']
        & V \arrow[d, "\phi_{\mathcal{B}_V}"] \\
        K^{(\kappa_1)} \arrow[r, "{\CoordMatrix[\mathcal{B}_U][\mathcal{B}_V]{T}}"']
                        \arrow[d, "{\CoordMatrix[\mathcal{B}_U][\mathcal{C}_U]{\Identity{U}}}"']
        & K^{(\kappa_2)} \arrow[d, "{\CoordMatrix[\mathcal{C}_V][\mathcal{B}_V]{\Identity{V}}}"] \\
        K^{(\kappa_1)} \arrow[r, "{\CoordMatrix[\mathcal{C}_U][\mathcal{C}_V]{T}}"']
        & K^{(\kappa_2)}
    \end{tikzcd}
    \]
    The two copies of $K^{(\kappa_1)}$ (resp.\ $K^{(\kappa_2)}$) carry different bases; the
    vertical maps are the matrices of the identity with respect to those two bases.
\end{remark}

\begin{remark}
    The matrices $\CoordMatrix[\mathcal{B}_V][\mathcal{C}_V]{\Identity{V}}$ and
    $\CoordMatrix[\mathcal{C}_V][\mathcal{B}_V]{\Identity{V}}$ are called
    \textbf{change of basis matrices} for $V$ and will be studied in detail in
    Section~\ref{sec:change-of-basis}.
\end{remark}

\subsection{Examples}

\begin{example}[Change of Bases Formula --- Endomorphism]
\label{ex:cob-endomorphism}
    Let $T: \bb{R}^2 \to \bb{R}^2$, $T(x,y) = (2x + y,\, x)$, with standard basis
    $\mathcal{B} = (e_0, e_1)$ and $\mathcal{C} = (e_0 + e_1,\, e_1)$.

    \textit{Step 1 --- $\CoordMatrix[\mathcal{B}][\mathcal{B}]{T}$}: Apply $T$ to $\mathcal{B}$:
    $T(e_0) = (2,1)$ and $T(e_1) = (1,0)$, so
    \[
        \CoordMatrix[\mathcal{B}][\mathcal{B}]{T}
        = \begin{pmatrix} 2 & 1 \\ 1 & 0 \end{pmatrix}
    \]

    \textit{Step 2 --- $\CoordMatrix[\mathcal{B}][\mathcal{C}]{\Identity{}}$}: Express
    $\mathcal{B}$ in $\mathcal{C}$: $e_0 = (e_0+e_1) - e_1$, $e_1 = e_1$, so
    \[
        \CoordMatrix[\mathcal{B}][\mathcal{C}]{\Identity{}}
        = \begin{pmatrix} 1 & 0 \\ -1 & 1 \end{pmatrix}
    \]

    \textit{Step 3 --- $\CoordMatrix[\mathcal{C}][\mathcal{B}]{\Identity{}}$}: Express
    $\mathcal{C}$ in $\mathcal{B}$: $e_0+e_1 = e_0+e_1$, $e_1 = e_1$, so
    \[
        \CoordMatrix[\mathcal{C}][\mathcal{B}]{\Identity{}}
        = \begin{pmatrix} 1 & 0 \\ 1 & 1 \end{pmatrix}
    \]

    \textit{Step 4 --- Apply the formula}:
    \[
        \CoordMatrix[\mathcal{C}][\mathcal{C}]{T}
        = \CoordMatrix[\mathcal{B}][\mathcal{C}]{\Identity{}}
          \cdot \CoordMatrix[\mathcal{B}][\mathcal{B}]{T}
          \cdot \CoordMatrix[\mathcal{C}][\mathcal{B}]{\Identity{}}
        = \begin{pmatrix} 1 & 0 \\ -1 & 1 \end{pmatrix}
          \begin{pmatrix} 2 & 1 \\ 1 & 0 \end{pmatrix}
          \begin{pmatrix} 1 & 0 \\ 1 & 1 \end{pmatrix}
        = \begin{pmatrix} 3 & 1 \\ 0 & -1 \end{pmatrix}
    \]

    \textit{Direct verification}: $T(e_0+e_1) = (3,1) = 3(e_0+e_1) + (-2)e_1$... let
    us instead verify via the coordinate formula.
    $\CoordVec[\mathcal{C}]{e_0+e_1} = \begin{psmallmatrix}1\\0\end{psmallmatrix}$,
    $\CoordVec[\mathcal{C}]{T(e_0+e_1)} = \CoordVec[\mathcal{C}]{(3,1)}$.
    Now $(3,1) = 3(e_0+e_1) - 2 e_1$, so
    $\CoordVec[\mathcal{C}]{(3,1)} = \begin{psmallmatrix}3\\-2\end{psmallmatrix}$
    --- but the first column of $\CoordMatrix[\mathcal{C}][\mathcal{C}]{T}$ is
    $\begin{psmallmatrix}3\\0\end{psmallmatrix}$. Let us recheck step 4:
    \[
        \begin{pmatrix} 1 & 0 \\ -1 & 1 \end{pmatrix}
        \begin{pmatrix} 2 & 1 \\ 1 & 0 \end{pmatrix}
        = \begin{pmatrix} 2 & 1 \\ -1 & -1 \end{pmatrix},
        \quad \text{then} \quad
        \begin{pmatrix} 2 & 1 \\ -1 & -1 \end{pmatrix}
        \begin{pmatrix} 1 & 0 \\ 1 & 1 \end{pmatrix}
        = \begin{pmatrix} 3 & 1 \\ -2 & -1 \end{pmatrix}
    \]
    So $\CoordMatrix[\mathcal{C}][\mathcal{C}]{T} = \begin{pmatrix} 3 & 1 \\ -2 & -1 \end{pmatrix}$.
    Direct check: first column is $\CoordVec[\mathcal{C}]{T(e_0+e_1)}
    = \begin{psmallmatrix}3\\-2\end{psmallmatrix}$ $\checkmark$; second column is
    $\CoordVec[\mathcal{C}]{T(e_1)} = \CoordVec[\mathcal{C}]{(1,0)}
    = \begin{psmallmatrix}1\\-1\end{psmallmatrix}$ $\checkmark$.
\end{example}

\begin{example}[Rank Computed from Any Representation]
    With $T$ and bases as in Example~\ref{ex:cob-endomorphism}:
    \[
        \Rank{\CoordMatrix[\mathcal{B}][\mathcal{B}]{T}}
        = \Rank{\begin{pmatrix} 2 & 1 \\ 1 & 0 \end{pmatrix}} = 2
        = \Rank{\begin{pmatrix} 3 & 1 \\ -2 & -1 \end{pmatrix}}
        = \Rank{\CoordMatrix[\mathcal{C}][\mathcal{C}]{T}}
    \]
    confirming Theorem~\ref{thm:rank-nullity-invariance}.
\end{example}
